\section{Experimental Setup}
\label{sec:experiments}

The experiments address four questions:
\begin{description}
  \item[RQ1:] Does typed-program supervision improve strict compositional execution accuracy
  over the same untuned base model and a strong cloud planner?
  \item[RQ2:] Does the specialised parser retain its accuracy under an independently generated
  surface form?
  \item[RQ3:] Which operator families and composition depths remain difficult?
  \item[RQ4:] Does the complete understand--plan--ground--execute--check runtime improve over its
  untuned counterpart on a balanced system snapshot?
\end{description}

\subsection{Frozen graph snapshot}

All reported executions use the same Parquet knowledge-graph snapshot and query-record
projection.  Table \ref{tab:dataset-stats} reports graph size and contract-level edge coverage.
Buyer coverage is \KGBuyerCoverage/\KGContracts (99.999\% after rounding), supplier coverage is
\KGSupplierCoverage/\KGContracts (94.87\%), category coverage is
\KGCategoryCoverage/\KGContracts (76.52\%), and at least one evidence edge is available for
\KGEvidenceCoverage/\KGContracts contracts (99.99\%).  These are structural coverage measures,
not correctness rates.

\begin{table}[t]
  \caption{Frozen graph and PACS evaluation statistics.  PACS family counts are intentionally
  non-uniform, so overall accuracy is a micro average and macro-family accuracy is also reported.}
  \label{tab:dataset-stats}
  \centering
  \small
  \begin{tabular}{llr}
    \toprule
    Block & Item & Count \\
    \midrule
    \multirow{4}{*}{Graph nodes}
      & Contract/notice & \KGContracts \\
      & Organisation & \KGOrgs \\
      & CPV & \KGCPV \\
      & Evidence & \KGEvidence \\
    \midrule
    \multirow{4}{*}{Contract coverage}
      & Buyer edge & \KGBuyerCoverage \\
      & Supplier edge & \KGSupplierCoverage \\
      & Category edge & \KGCategoryCoverage \\
      & At least one evidence edge & \KGEvidenceCoverage \\
    \midrule
    \multirow{6}{*}{PACS}
      & Each surface channel & \PacsN \\
      & Answerable / unsupported / no result & 694 / 138 / 90 \\
      & Level L1 / L2 / L3 & 279 / 339 / 304 \\
      & Seen / declared unseen shapes & 431 / 491 \\
      & Smallest family (F7) & 67 \\
      & Largest families (F2 and F6) & 149 each \\
    \bottomrule
  \end{tabular}
\end{table}

\subsection{PACS construction}
\label{sec:pacs-construction}

PACS is designed to test combinations of closed-algebra primitives rather than broad world
knowledge.  Seven families cover: F1 spending and volume; F2 temporal change; F3 supplier
concentration; F4 cross-buyer comparison; F5 overlap and exclusion; F6 relational composition;
and F7 disclosure/compliance.  L1 questions use an atomic or shallow computation, L2 questions
combine two or three operations, and L3 questions use nested set membership, filtered
comparisons, set algebra, or group composition.

Construction proceeds as follows.  First, a family template samples anchors from the frozen
record universe and instantiates a gold tree.  Second, the tree must pass static validation,
execute successfully in both implementations, produce agreeing answers, and satisfy a
family-specific non-degeneracy rule.  Third, supported variants are paired with out-of-grammar
requests and scopes that produce a clean empty denotation.  Fourth, the generator records a
tree hash and a shape signature that removes literals while preserving operator structure.
Fifth, whole intent clusters---including their status variants---are assigned to development or
test with a fixed seed, so closely related variants do not cross the split.

Each test intent has two evaluated surfaces.  Channel A is a deterministic independent rendering
of the semantics.  Channel B is a separately naturalised question that must preserve required
literals and pass a logic gate.  The two channels share the gold program and oracle, allowing a
paired robustness comparison without changing the underlying computation.

The sealing code enforces exact normalised question non-overlap with the training export, gold
tree-hash non-overlap, and absence of each row labelled \emph{unseen} from the set of training
shape signatures.  It also screens Channel B for trigram Jaccard similarity of at least 0.7
against every seventh training trigram set.  The source documentation names five gates, but no
separate G3 test is implemented and the near-duplicate screen is subsampled.  We therefore call
PACS a \emph{mechanically screened compositional test}, not an exhaustive zero-leakage or
maximum-compound-divergence split.  The test-channel SHA-256 hashes and construction commands are
listed in Appendix \ref{app:reproducibility}.

\subsection{Systems}

The primary PACS comparison uses the same algebra description and one-question prompt for all
systems:
\begin{itemize}
  \item \textbf{Untuned 8B}: Qwen3-8B emits a tree or abstention without task-specific
  fine-tuning.
  \item \textbf{Cloud planner}: \texttt{grok-4-1-fast-non-reasoning} uses the same free-decoding
  interface and algebra contract.
  \item \textbf{Typed-program 8B}: the Qwen3-8B QLoRA adapter described in Section
  \ref{sec:program-supervision}.
\end{itemize}
The main protocol uses temperature zero, a single free decode, no guided JSON grammar, and no
reflection/repair.  This preserves comparability with the cloud model.  Only the specialised
parser was evaluated on Channel B; Channel A therefore provides the fair three-system
comparison, while A--B measures within-model surface robustness.

The secondary end-to-end comparison contains 13 categories with 20 cases each: ambiguous,
no-result, and unsupported abstention; Boolean, bridge-join, categorical, comparison, count,
factoid, minimum/maximum, set, sum, and top-$k$ questions.  We compare the full local runtime,
the same Qwen3 stack without specialised adaptation, and the stored cloud-teacher run.  These
results are tied to frozen summaries under the external \texttt{outputs} artifact store; the
exact model/run manifest is less complete than for PACS, so they are reported as secondary
system evidence rather than the main benchmark.

\subsection{Metrics and statistical reporting}

The primary metric is strict accuracy from Equation \ref{eq:strict-correctness}.  It credits an
answerable item only when a valid tree executes and its typed answer matches the oracle; an
unsupported item requires explicit abstention.  The supplementary \emph{safe accuracy} also
credits a literal-faithful multi-answer or empty-result execution for eligible non-answerable
cases.  Because this metric mixes exact answers and operationally safe outcomes, it is relegated
to Appendix \ref{app:full-results}.  Tree-valid rate measures only static validity and is not a
correctness metric.

We report exact numerators and denominators.  Parenthesised intervals in the main tables are 95\%
Wilson score intervals for a binomial proportion; they describe uncertainty on the fixed test
items, not model-training variation.  All headline model evaluations are single deterministic
decodes, so they do not estimate sensitivity to checkpoint, prompt, provider, or hardware
variation.  Family macro accuracy is the unweighted mean across F1--F7; overall accuracy is the
micro average across all \PacsN{} cases.
